\newcommand{\docshorttitle}{SUPPL.\ MEMO P\&A --- \S 425.17(c) COMMERCIAL CARVE-OUT (CGC-25-631802)}
\newcommand{\doctitleblock}{PLAINTIFF'S SUPPLEMENTAL MEMORANDUM OF POINTS AND AUTHORITIES IN FURTHER OPPOSITION TO DEFENDANT'S SPECIAL MOTION TO STRIKE (CODE CIV.\ PROC., \S\S~425.16, 425.17(c))}
\newcommand{\docsubblock}{(Cal.\ Rules of Court, rules 3.1113, 3.1306(c); supplemental to Plaintiff's previously-filed Opposition; cross-noticed for hearings of May 13, 2026 and May 19, 2026, Dept.\ 302)}
% Pleading-paper preamble for APRIL-24-2026 cover sheets and oral-argument
% outlines (CGC-25-631801 and CGC-25-631802). Matches house style from
% APRIL-23-SAC-DEFENSE/court-pdfs/REPLY-SAC-PREAMBLE.tex, simplified for
% single-to-three-page artifacts that do not require TOC/TOA infrastructure.
%
% Cross-hearing caption (CAPTION-801/802.tex): use \CaptionDocTitleBlock,
% \CaptionDocSubBlock, and the crosshearingsschedule environment (defined below)
% so long \doctitleblock text and hearing dates stay inside the 3in minipage.
\documentclass{article}
\usepackage[utf8]{inputenc}
\usepackage{graphicx}
\usepackage{geometry}
    \geometry{top=1in, bottom=2.2in, left=1.0in, right=1.0in, textheight=336pt}
    \setlength{\footskip}{70pt}
\usepackage{ulem}
\usepackage{fancyhdr}
    \renewcommand{\headrulewidth}{0pt}
    \renewcommand{\footrulewidth}{0.4pt}
\usepackage{parskip}
    \setlength{\parskip}{0mm}
\usepackage{quoting}
    \quotingsetup{leftmargin=0.5in, rightmargin=0.5in, vskip=-1.5mm}
    \makeatletter
        \g@addto@macro\quoting\singlespacing
        \g@addto@macro\quoting{\vspace{-2mm}}
    \makeatother
\usepackage{mathptmx}
\usepackage{amssymb}
\renewcommand{\rmdefault}{ptm}
\renewcommand{\normalsize}{\fontsize{12}{12}\selectfont}
\newcommand*{\ptm}{\fontfamily{ptm} \selectfont \fontsize{12}{12}\selectfont}
\usepackage{setspace}
    \doublespacing
\newcommand{\tab}{\hspace*{.0in}}
\newenvironment{tightcenter}{%
    \setlength\topsep{0pt}
    \setlength\parskip{0pt}
    \begin{center}
}{%
    \end{center}
}
\usepackage{tikz}
\usetikzlibrary{calc}
% Pleading-paper vertical double rule at the left margin.
\AddToHook{shipout/background}{%
    \begin{tikzpicture}[remember picture, overlay]
        \draw[line width=0.5pt] ($(current page.north west) + (0.75in,0)$) -- ($(current page.south west) + (0.75in,0)$);
        \draw[line width=0.5pt] ($(current page.north west) + (0.80in,0)$) -- ($(current page.south west) + (0.80in,0)$);
    \end{tikzpicture}%
}
\usepackage[left, pagewise]{lineno}
\setlength{\linenumbersep}{0.35in}
\renewcommand{\linenumberfont}{\normalfont\small}
\usepackage{titlesec}
\usepackage[hidelinks]{hyperref}
\usepackage{bookmark}
\usepackage{textcomp}
\usepackage{longtable}
\usepackage{booktabs}
\usepackage{array}
\usepackage{multirow}
\usepackage{calc}
% --- Cross-hearing caption cover (3in minipage): compact title + 3-column hearing table ---
% Use on first page only; keeps long \doctitleblock / \docsubblock on one sheet with party caption.
\newcommand{\CaptionDocTitleBlock}{%
  \begingroup
  \fontsize{10}{11.5}\selectfont
  \bfseries
  \raggedright
  \setlength{\parskip}{0pt}%
  \noindent\doctitleblock\par
  \endgroup
}
\newcommand{\CaptionDocSubBlock}{%
  \begingroup
  \footnotesize
  \raggedright
  \setlength{\parskip}{0pt}%
  \noindent\docsubblock\par
  \endgroup
}
% #1 = department number only (e.g.\ 301 or 302) for the line ``Dept.\ #1''.
\newenvironment{crosshearingsschedule}[1]{%
  \par\begingroup
  \footnotesize
  \setlength{\tabcolsep}{2.2pt}%
  \renewcommand{\arraystretch}{1.04}%
  \setlength{\parskip}{0pt}%
  \setlength{\parindent}{0pt}%
  \noindent\textbf{Cross-noticed for pending defense hearings (this case, Dept.\ #1):}\par\vspace{2pt}%
  \begin{tabular}{@{}>{\raggedright\arraybackslash}p{0.11\linewidth}>{\raggedright\arraybackslash}p{0.34\linewidth}>{\raggedright\arraybackslash}p{0.53\linewidth}@{}}%
  \textbf{\#} & \textbf{Motion} & \textbf{Date / dept / judge} \\ \hline
}{%
  \end{tabular}\par\endgroup\smallskip
}
% Pandoc pipe-table column widths use \real{#}.
\newcommand{\real}[1]{\pgfmathparse{#1}\pgfmathresult}
\titleformat{\section}{\fontsize{13}{13}\selectfont\normalfont\bfseries\uppercase}{}{0em}{}
\titleformat{\subsection}{\fontsize{12}{12}\selectfont\normalfont\bfseries}{}{0em}{}
\titleformat{\subsubsection}{\normalfont\bfseries}{}{0em}{}
\titlespacing*{\section}{0pt}{6pt}{2pt}
\titlespacing*{\subsection}{0pt}{4pt}{2pt}
\titlespacing*{\subsubsection}{0pt}{2pt}{2pt}
\newcommand{\calrules}[1]{Cal. Rules of Court, rule #1}
\newcommand{\ccp}[1]{Code Civ. Proc., \textsection~#1}
\newcommand{\evidcode}[1]{Evid. Code, \textsection~#1}
\providecommand{\tightlist}{%
  \setlength{\itemsep}{0pt}\setlength{\parskip}{0pt}}
% Footer: page N of total + document short-title.
\usepackage{lastpage}
\pagestyle{fancy}
\fancyhf{}
\fancyfoot[C]{\thepage{} of \pageref{LastPage} \\[1pt] \footnotesize \docshorttitle}
\fancyfoot[R]{}

\begin{document}
% Cross-hearing caption for CGC-25-631802 (Dept. 302, Hon. Quinn).
\nolinenumbers
\begin{singlespace}
\noindent Franciscus Dylan Rosario \\
7624 Melody Drive \\
Rohnert Park, CA 94928 \\
Tel: 650.488.7510 \\
E: soltrinox@gmail.com \\
Plaintiff, In Pro Per
\end{singlespace}
\vspace*{9mm}
\begin{tightcenter}
\textbf{SUPERIOR COURT OF THE STATE OF CALIFORNIA} \\
\textbf{COUNTY OF SAN FRANCISCO}
\end{tightcenter}
\vspace*{3mm}
\begin{minipage}[t]{3in}
    \raggedright
    \textbf{FRANCISCUS DYLAN ROSARIO,} \\ \textbf{PLAINTIFF}, \\
    \hspace{1cm} v. \\
    \small
    CSAA Insurance Exchange; and Does 1 through 50, inclusive;\\
    \normalsize
    \textbf{DEFENDANTS}
    \makebox[3in]{\hrulefill}
\end{minipage}%
\hspace{3mm}%
\begin{minipage}[t]{0.5pt}
    \vspace{0pt}
    \rule{0.5pt}{2.42in}
\end{minipage}%
\hspace{3mm}%
\begin{minipage}[t]{3in}
    \raggedright
    \noindent\textbf{Case No.: CGC-25-631802}\par\smallskip
    \CaptionDocTitleBlock\smallskip
    \CaptionDocSubBlock\smallskip
    \begin{singlespace}
    \setlength{\parindent}{0pt}
    \begin{crosshearingsschedule}{302}
    (1) & MFL / leave (SAC) & \textbf{Apr.\ 30, 2026} \\
    (2) & Anti-SLAPP (\ccp{425.16}) & \textbf{May 13, 2026, 9:30 a.m.};\ Dept.\ 302;\ Hon.\ Joseph M.\ Quinn \\
    (3) & Strike (\ccp{436}) & \textbf{May 19, 2026};\ Dept.\ 302;\ Hon.\ Joseph M.\ Quinn \\
    (4) & Demurrer & Hrng.\ date \textbf{[TBD]} (reserved) \\
    \end{crosshearingsschedule}
    \noindent Complaint filed: Dec.\ 23, 2025 \\
    \noindent FAC filed: Feb.\ 25, 2026 \\
    \end{singlespace}
\end{minipage}
\vspace*{8mm}
\newpage
\linenumbers


\section*{I. INTRODUCTION AND VEHICLE}

This Supplemental Memorandum of Points and Authorities is respectfully submitted in further opposition to Defendant CSAA Insurance Exchange's special motion to strike under Code of Civil Procedure section 425.16. It is offered as a supplement to --- and not as a substitute for --- Plaintiff's previously-filed Opposition. The vehicle is \calrules{3.1113} (memorandum of points and authorities) read together with \calrules{3.1306(c)} (lodging of supplemental authority for hearings), as the supporting Notice of Lodging Supplemental Authority, the Supplemental Request for Judicial Notice, and the bench briefs filed concurrently in this same cross-hearing package make clear. Plaintiff incorporates the previously-filed Opposition by reference and submits this supplemental memorandum to focus the Court's attention on the structurally prior question raised by Code of Civil Procedure section 425.17(c): whether the special motion to strike procedure is available to Defendant at all.

As shown below, the answer is no. CSAA is, by name, the paradigmatic Section 425.17(c) defendant; its pre-litigation claim-handling communications to Plaintiff are the paradigmatic Section 425.17(c) conduct; and the audience for those communications --- Plaintiff himself, an actual customer of the regulated insurance-claims-handling product --- is the paradigmatic Section 425.17(c) audience. Because the carve-out applies, \ccp{425.16} does not apply to Plaintiff's causes of action, and the special motion to strike must be denied at the gateway without reaching the merits-stage \ccp{425.16(b)(1)} ``probability of prevailing'' inquiry.

The Supplemental Request for Judicial Notice filed concurrently reproduces the verbatim text of \ccp{425.17}, \ccp{425.16(e), (g), (i)}, \ccp{904.1(a)(13)}, Civil Code \S~47(b), Business and Professions Code \S~6128(a), Insurance Code \S~790.03(h), and 10 California Code of Regulations \S~2695.1. A companion Supplemental Memorandum filed concurrently in this case (the ``Customer-Doctrine Supplemental'') addresses the discrete question whether a third-party claimant is a ``customer'' within the meaning of section 425.17(c)(2); that question is treated only summarily here.

\section*{II. STATEMENT OF FACTS GERMANE TO THE \S~425.17(c) GATEWAY}

The sole moving Defendant in CGC-25-631802 is CSAA Insurance Exchange, a California-licensed insurer. Plaintiff's pleaded causes of action arise from CSAA's pre-litigation claims-handling conduct, including: (a) communications made to Plaintiff as the claimant-customer; (b) representations concerning the status and conclusions of the claim investigation; and (c) the issuance of a denial communication. Each of those communications was made in the course of CSAA's delivery of the insurance services for which premiums had been paid. The intended audience throughout was Plaintiff himself --- an actual customer of CSAA. The conduct also took place within the context of the regulatory regime administered by the California Department of Insurance under \S~790.03(h) and the Fair Claims Settlement Practices Regulations at 10 CCR \S~2695.1 et seq.

\section*{III. LEGAL STANDARD --- \S~425.17(c) IS A STRUCTURAL GATEWAY}

\subsection*{A. The Carve-Out Is Categorical and Sequentially Prior to \S~425.16.}

The Legislature enacted Code of Civil Procedure section 425.17 in 2003 because it determined that \ccp{425.16} was being misused. \ccp{425.17(a)} declares: ``[T]here has been a disturbing abuse of Section 425.16, the California Anti-SLAPP Law, which has undermined the exercise of the constitutional rights of freedom of speech and petition for the redress of grievances, contrary to the purpose and intent of Section 425.16.'' Subdivision (c) is the Legislature's response to one identifiable category of abuse: special motions to strike filed by commercial actors --- including, by name, those engaged in the business of insurance --- to defeat causes of action that arise from their own commercial representations or from conduct undertaken in the course of delivering their goods or services to actual or potential customers.

The California Supreme Court has confirmed in \textit{Simpson Strong-Tie Co.\ v.\ Gore} (2010) 49 Cal.4th 12, and the Court of Appeal has reaffirmed in \textit{JAMS, Inc.\ v.\ Superior Court} (2016) 1 Cal.App.5th 984 and \textit{Stewart v.\ Rolling Stone LLC} (2010) 181 Cal.App.4th 664, that the carve-out is categorical: when its elements are met, \ccp{425.16} ``does not apply'' at all. The Court does not weigh competing interests; it declines to extend the special motion to strike procedure to the category of disputes the Legislature has removed from that mechanism. \textit{Club Members for an Honest Election v.\ Sierra Club} (2008) 45 Cal.4th 309 supplies the sequencing rule: where \ccp{425.17} is invoked, the trial court must address the carve-out before any \ccp{425.16} analysis. This sequencing is not optional.

\subsection*{B. The Two-Element Test of \S~425.17(c).}

The two elements of \ccp{425.17(c)} are straightforward. \textbf{First}, the defendant must be ``primarily engaged in the business of selling or leasing goods or services,'' with insurance specifically enumerated. \textbf{Second}, the cause of action must arise from a statement or conduct that consists either of representations of fact about the defendant's business operations, goods, or services, made for the purpose of obtaining approval for, promoting, or securing a commercial transaction, \textbf{or} of conduct made in the course of delivering those goods or services. In addition, \ccp{425.17(c)(2)} requires that the intended audience be (i) an actual or potential buyer or customer; (ii) a person likely to repeat the statement to or otherwise influence such a customer; or (iii) alternatively, that the conduct have arisen out of or within the context of a regulatory approval process, proceeding, or investigation. Each of these audience branches is independently sufficient.

\section*{IV. ARGUMENT --- THE \S~425.17(c) CARVE-OUT REQUIRES DENIAL OF THE SPECIAL MOTION TO STRIKE}

\subsection*{A. CSAA Is the Paradigmatic \S~425.17(c) Defendant.}

CSAA is a California-licensed insurer. Insurance is one of the named industries in the text of \ccp{425.17(c)}. The first element is satisfied as a matter of law.

\subsection*{B. Plaintiff's Causes of Action Arise from Conduct ``Made in the Course of Delivering'' Insurance Services.}

Each pre-litigation claim-handling communication --- representations about the status of investigation, communications regarding the basis for denial, and the denial communication itself --- was, by its nature, ``made in the course of delivering'' the insurance services for which premiums had been paid. That is the very phrase \ccp{425.17(c)(1)} uses to capture commercial conduct that cannot be re-characterized as protected speech for anti-SLAPP purposes.

\subsection*{C. The Intended Audience Was Plaintiff, an Actual Customer of the Insurer.}

The intended audience for each operative communication was Plaintiff himself --- an actual customer of CSAA. That alone satisfies \ccp{425.17(c)(2)}. Independently, to the extent any communication can be said to have ``arisen out of or within the context of'' the Department of Insurance's regulatory regime --- which it did, because Insurance Code \S~790.03(h) and 10 CCR \S~2695.1 et seq. directly govern misrepresentations to claimants, the duty to adopt reasonable investigation standards, and the obligation to attempt prompt, fair, and equitable settlement --- the alternative branch of \ccp{425.17(c)(2)} is independently triggered.

\subsection*{D. Defendant's ``Free Speech'' Re-Characterization Fails.}

Defendant resists this conclusion by re-characterizing the operative conduct as ``petitioning'' or ``speech in connection with an issue of public interest'' under \ccp{425.16(e)}. That re-characterization fails for three independent reasons.

\textbf{First}, the conduct at issue is not speech directed at, or in connection with, any official proceeding; it is commercial speech directed at the insured during the adjustment of a private contractual claim. The mere fact that a denial letter eventually became a subject of litigation does not retroactively convert pre-litigation claim handling into ``petitioning activity.'' \textit{Rusheen v.\ Cohen} (2006) 37 Cal.4th 1048, and the cases following it, draw a clear line between communicative acts integral to a judicial proceeding and the underlying non-communicative or pre-litigation conduct that gives rise to a cause of action. The gravamen of Plaintiff's claims is the latter.

\textbf{Second}, even where some element of the conduct could plausibly be styled as speech, \ccp{425.17(c)} governs by its own terms; the carve-out is not displaced by the breadth of \ccp{425.16(e)}. The Legislature wrote \ccp{425.17} precisely to override that kind of categorical re-characterization in the commercial context.

\textbf{Third}, \textit{Flatley v.\ Mauro} (2006) 39 Cal.4th 299, reaffirmed in \textit{Lefebvre v.\ Lefebvre} (2011) 199 Cal.App.4th 696, holds that conduct conceded or shown as a matter of law to be illegal is not protected by \ccp{425.16} in the first instance. To the extent Defendant's special motion depends on conduct that, on an uncontroverted record, violates \S~790.03(h) or constitutes extrinsic fraud, the \textit{Flatley} gateway closes regardless of how the conduct is labeled.

\subsection*{E. The ``Concealment of the Failure to Investigate'' Theory Is Squarely Within \S~425.17(c)(1).}

Plaintiff's pleaded theory --- that the insurer communicated a denial premised on a represented conclusion of investigation when in fact no adequate investigation had occurred --- is structurally a \S~790.03(h)(1) and (h)(3) theory dressed in common-law clothing. The Court need not decide on this motion whether each statutory element has been violated; it need only recognize that conduct alleged to have violated \S~790.03(h) is, by definition, conduct ``in the course of delivering'' insurance services and therefore squarely within \ccp{425.17(c)(1)}. Whether Plaintiff ultimately prevails on the merits is not the question on a \ccp{425.16} special motion; the question is whether Defendant may use \ccp{425.16} to extract that theory from the ordinary civil process. \ccp{425.17(c)} answers that question in the negative.

\subsection*{F. The Cross-Noticed \S~436 Motion Does Not Rescue the Defense.}

Defendant has cross-noticed a motion under \ccp{436} that incorporates by reference the ``everything is petitioning'' rhetoric of its anti-SLAPP papers. \textit{PH II, Inc.\ v.\ Superior Court} (1995) 33 Cal.App.4th 1680, 1682--83, holds that a \ccp{436} motion is not a backdoor demurrer; by parity of reasoning, it is not a backdoor anti-SLAPP motion either. To the extent the \ccp{436} papers attempt to import \ccp{425.16(e)} themes, they collide with the same \ccp{425.17(c)} carve-out that defeats the special motion: the underlying conduct is commercial-insurer claims handling, not protected petitioning.

\subsection*{G. \S~425.17(e) --- No Interlocutory Appeal, No Stay, No Writ.}

If this Court denies the special motion to strike on \ccp{425.17} grounds --- as it must on this record --- the denial is not subject to interlocutory appeal under \ccp{904.1(a)(13)}. \ccp{425.17(e)} is explicit. The Legislature concluded that the ordinary anti-SLAPP appeal route is itself a vehicle for abuse where the underlying claim was never within \ccp{425.16} to begin with. The Court of Appeal cannot be drafted into delay where the trial court has done what \ccp{425.17} directs: applied the carve-out, declined the special motion, and returned the matter to the ordinary civil docket.

\section*{V. CONCLUSION}

CSAA is the paradigmatic \ccp{425.17(c)} defendant. Its pre-litigation communications to Plaintiff are the paradigmatic \ccp{425.17(c)} conduct. The audience for those communications was the paradigmatic \ccp{425.17(c)} audience. There is no jurisprudential or textual route by which Defendant can carry its threshold burden to invoke \ccp{425.16}. Plaintiff respectfully requests that the Court deny the special motion to strike at the gateway under \ccp{425.17(c)}, without reaching the second prong of the anti-SLAPP analysis, and proceed with the ordinary civil docket in CGC-25-631802.

\input{SIGBLOCK.tex}
\end{document}
